\section{A finite certificate for an infinite matrix}\label{cq:certificate}

This appendix specifies the rational calculation cited after
Theorem~\ref{cq:twins}. The accompanying file
\texttt{code/verify\_twin\_certificate.py} uses integer and rational
arithmetic for every decisive comparison. Its decimal output is only
a readable summary.

Let \(Y=100\), and write
\[
 A_Y(z)=\prod_{2<p\leq Y}\frac{1-z/(p+3)}{1-z/p}
       =\sum_{n\geq0}\alpha_nz^n,\qquad
 \frac{F(1-z)}{F(1)}=A_Y(z)B_Y(z).
\]
The tail estimate in the main text gives
\[
 B_Y(z)=\exp\left(\sum_{k\geq1}\frac{\epsilon_kz^k}{k}\right)
       =\sum_{m\geq0}\beta_mz^m,\qquad
 0\leq\epsilon_k\leq3Y^{-k}.
\]
Coefficient comparison with \((1-z/Y)^{-3}\) yields
\begin{equation}\label{cq:betabound}
      \beta_0=1,\qquad 0\leq\beta_m\leq\binom{m+2}{2}Y^{-m}.
\end{equation}
Keeping one common sequence \((\beta_m)\) preserves the dependence between
the matrix entries; independent intervals for the entries would throw away
useful structure.

Put \(S=\{3,5,7,11,13,17,19,23\}\) and \(T=\{5,7,13,19\}\). For \(p\in T\)
define the degree-fourteen polynomial
\[
 P_p(u)=
 \frac{u^7\prod_{q\in S\setminus\{p\}}(u-1/q)}
      {p^{-7}\prod_{q\in S\setminus\{p\}}(1/p-1/q)}.
\]
It has rational coefficients and satisfies
\(P_p(1/q)=\mathbf1_{\{p=q\}}\) for \(q\in T\).
Let \(P_i(u)=\sum_r a_{i,r}u^r\), and form
\[
 (Q_m)_{ij}=\sum_{r,s}a_{i,r}a_{j,s}\alpha_{r+s+2-m},
 \qquad \alpha_n=0\quad(n<0).
\]
The actual infinite moment form on these four coefficient vectors is
\[
                         \sum_{m=0}^{30}\beta_m Q_m.
\]
All \(\alpha_n\) through degree \(30\) are obtained from the rational
recursion
\[
 \alpha_0=1,\qquad
 \alpha_n=\frac1n\sum_{k=1}^n
 \left(\sum_{2<p\leq100}\{p^{-k}-(p+3)^{-k}\}\right)\alpha_{n-k}.
\]
Thus the certificate has no unspecified numerical inputs.

Exact symmetric elimination verifies that \(Q_0\) is negative definite.
Rational Gram--Schmidt changes the four polynomials to a basis in which
\(Q_0=-\operatorname{diag}(a_1,a_2,a_3,a_4)\), all \(a_i>0\).
In this basis, exact elimination also verifies that \(Q_1,Q_2\) are
negative definite. Define
\[
 E_{ij}=\sum_{m=3}^{30}\binom{m+2}{2}100^{-m}|(Q_m)_{ij}|.
\]
The four exact comparisons checked by the program are
\begin{equation}\label{cq:rowcertificate}
                    7\sum_jE_{ij}<a_i\qquad(1\leq i\leq4).
\end{equation}
For orientation, the four ratios \(\sum_jE_{ij}/a_i\), rounded only for
display, are
\[
        0.001350780,\quad0.003829068,\quad
        0.027644920,\quad0.140852054 .
\]

For any nonzero real vector \(v\), use \eqref{cq:betabound}, discard the
two negative tail contributions, and apply
\(2|v_i v_j|\leq v_i^2+v_j^2\). The result is
\[
 v^{\mathsf T}\left(\sum_{m=0}^{30}\beta_mQ_m\right)v
 \leq-\sum_i\left(a_i-\sum_jE_{ij}\right)v_i^2
 <-\frac67\sum_i a_i v_i^2<0 .
\]
The four polynomials are independent by their interpolation values.
Hence they certify four negative directions in the infinite matrix.

The certificate's witnesses use the already known twin upper endpoints
\(5,7,13,19\). Its purpose is to show how a statement about infinitely
many uncomputed primes can be certified by finite exact inequalities.
It makes no claim to discover a new twin pair.
